\documentclass[12pt,a4paper]{article}

\usepackage[T1]{fontenc}
\usepackage[utf8]{inputenc}
\usepackage[romanian]{babel}
\usepackage{lmodern}
\usepackage{microtype}

\usepackage{hyperref}
\hypersetup{
    colorlinks=true,
    urlcolor=blue,
    linkcolor=black
}

\setlength{\parskip}{0.8em}
\setlength{\parindent}{0pt}

\usepackage{titling}

\pretitle{\begin{center}\LARGE\bfseries\vspace{-1cm}}
\posttitle{\par\vspace{0.5em}\hrule\vspace{0.8em}\end{center}}
\preauthor{\begin{center}\large}
\postauthor{\end{center}}
\predate{\begin{center}\small}
\postdate{\end{center}}

\title{Pacient virtual stomatologic}
\author{
  \textsc{Marcu Emanuel}\\[0.4em]
  \normalsize\textbf{Echipa: VirtuDent}
}
\date{}

\begin{document}

\maketitle
\vspace{1em}

\section{Funcționalități de bază ale aplicației}

Principalele funcționalități sunt:
\begin{itemize}
    \item Simulare realistă de pacient, care descrie simptomele folosind un limbaj informal, asemănător cu cel al unui pacient real;
    \item Conversație liberă între student și pacient, în care întrebările formulate influențează răspunsurile primite;
    \item Generarea automată a răspunsurilor în limbaj natural;
    \item Oferirea de hint-uri progresive, care îl ajută pe student să formuleze întrebări mai relevante, fără a dezvălui direct diagnosticul;
    \item Posibilitatea de a solicita investigații suplimentare, precum o radiografie, preluată dintr-o bază de date cu imagini specifice fiecărei afecțiuni;
    \item Formularea și evaluarea automată a diagnosticului, cu feedback detaliat privind corectitudinea și raționamentul;
    \item Salvarea istoricului conversațiilor și rezultatelor, pentru urmărirea progresului fiecărui utilizator.
\end{itemize}

\section{Descrierea problemei rezolvate cu ajutorul AI}

Aplicația propusă urmărește să sprijine studenții la medicină dentară în procesul de formare clinică printr-un pacient virtual interactiv.  
Problema identificată este faptul că studenții intră în contact direct cu pacienții abia în anul IV, iar lipsa experienței practice duce la ezitare și nesiguranță în stabilirea unui diagnostic.

Prin acest proiect se oferă un mod sigur și realist de a exersa anamneza și raționamentul clinic, simulând un pacient care își descrie simptomele în limbaj informal.  
Studentul poate adresa întrebări, primi răspunsuri, cere indicii (hinturi) și, în final, propune un diagnostic care este verificat automat.

\noindent\textbf{Utilizatori:} studenți la medicină dentară și cadre didactice.\\
\textbf{Date de intrare:} întrebările formulate de student și diagnosticul propus (text).\\
\textbf{Date de ieșire:} răspunsurile pacientului, hinturile și evaluarea diagnosticului.\\
\textbf{Tipul datelor:} text în limbaj natural, ulterior posibil și imagini radiologice.

\medskip
Performanța sistemului AI este evaluată prin calitatea simulării pacientului virtual și acuratețea evaluării diagnosticului propus de student.  
Se urmărește ca răspunsurile generate să fie realiste, logice și consecvente cu simptomele definite, iar verdictul final („corect”, „parțial”, „greșit”) să reflecte cât mai fidel situația cazului.  
Astfel, performanța AI-ului indică capacitatea sistemului de a reproduce comportamentul unui pacient real și de a evalua obiectiv răspunsurile studentului.

\section{Related Work \& Useful Tools and Technologies}

Pentru dezvoltarea aplicației au fost analizate proiecte similare.

\subsection{Virtual Patient Simulator for Skill Training in Dentistry – University of Peradeniya, Sri Lanka}
\href{https://github.com/cepdnaclk/e16-4yp-Virtual-Patient-Simulator-for-Skill-Training-in-Dentistry}{Link GitHub}

Proiectul propune un simulator educațional care combină realitatea virtuală (VR) și tehnici de inteligență artificială pentru antrenarea abilităților clinice ale studenților la medicină dentară.

\textbf{Date utilizate:}
Două cazuri clinice virtuale, fiecare conținând informații despre anamneză, examinare, investigații și diagnostic.  
Au fost incluse modele 3D intraorale realizate în Blender și imagini radiologice (DPT, IOPA, CBCT).

\textbf{Algoritmi de AI folosiți:}
Sistemul utilizează un chatbot bazat pe reguli predefinite, care oferă răspunsuri și feedback automat pentru student.  
Nu sunt folosite modele NLP sau algoritmi de învățare automată avansați, ci o logică deterministă bazată pe scenarii.

\textbf{Performanțe obținute:}
Simulatorul a fost testat pe un grup de 33 de studenți, împărțiți între metoda tradițională și cea virtuală.  
Scorul mediu la testul OSCE a crescut de la 44.21 la 58.95 pentru grupul VP.  
Rezultatele sunt similare cu cele ale grupului de control (p-value = 0.818 pre-test, 0.666 post-test).  
Concluzia studiului: performanțele sistemului VP sunt comparabile cu instruirea clinică reală.  
Studenții au raportat o mai bună încredere în diagnostic și o motivație crescută pentru exersare.

\textbf{Librării și tehnologii:}
React, Redux, Firebase Firestore pentru interfața web și gestiunea datelor; Blender pentru modelarea 3D.

\textbf{Git / resurse:}
Codul sursă și documentația completă sunt disponibile public pe GitHub.

\subsection{Deep Learning-Based Chatbot for Medical Assistance – Puneet Shivaay (GitHub)}
\href{https://github.com/PuneetShivaay/Deep-Learning-Based-Chatbot-For-Medical-Assistance}{Link GitHub}

Proiectul dezvoltă un chatbot inteligent pentru asistență medicală, care oferă recomandări de diagnostic pe baza simptomelor introduse de utilizator.  
Scopul este de a crea o interacțiune simplă, în limbaj natural, între pacient și un asistent virtual.

\textbf{Date utilizate:}
Modelul este antrenat pe un fișier JSON (\textit{intents.json}), ce conține perechi între simptome și afecțiuni posibile, precum și mai multe variante de formulare a întrebărilor.  
În plus, sunt salvate listele de cuvinte și etichete în fișierele \textit{words.pkl} și \textit{labels.pkl}, pentru a facilita recunoașterea tiparelor de text.

\textbf{Algoritmi de AI folosiți:}
Chatbotul folosește o rețea neuronală antrenată cu \textbf{TensorFlow/Keras}, împreună cu tehnici de procesare a limbajului natural (NLP).  
Textul introdus este transformat într-un vector (\textit{bag-of-words}), iar modelul învață să asocieze întrebările utilizatorului cu intențiile și diagnosticele corespunzătoare.

\textbf{Performanțe obținute:}
Proiectul are un caracter demonstrativ, fără metrici numerice (precizie, recall etc.), însă modelul reușește să ofere răspunsuri logice și coerente pentru majoritatea întrebărilor din setul de test.  
Scopul este validarea ideii și testarea funcționării chatbotului.

\textbf{Librării și tehnologii:}
Aplicația este realizată în \textbf{Python 3.7}, folosind TensorFlow/Keras pentru modelul AI, NLTK pentru procesarea textului, Tkinter pentru interfața grafică și Pickle pentru salvarea datelor de antrenament.

\textbf{Git / resurse:}
Codul sursă, modelul antrenat (\textit{chatbot\_model.h5}) și fișierele auxiliare sunt disponibile public pe GitHub.

\end{document}
